\documentclass[10pt,a4paper]{amsart}
\usepackage[margin=19mm]{geometry}
\usepackage{amsmath,amssymb,mathtools}
\usepackage[scr=rsfs,cal=euler]{mathalfa}
\usepackage{xfrac}
\usepackage[unicode,hidelinks,bookmarks=false]{hyperref}
\newcommand{\ie}{i.e.\ }
\newcommand{\cf}{cf.\ }
\newcommand{\resp}{respectively\ }
\newcommand{\Subsection}{\subsection}
\providecommand{\nobreakdash}{\nobreak\hbox{-}}
\setlength{\emergencystretch}{2em}
\multlinegap0pt
\usepackage{bm}
\usepackage{bbm}
\usepackage{dsfont}
\usepackage{xspace}
\usepackage{cancel}
\usepackage{mathtools}
\usepackage{cleveref}
\usepackage{amsthm}
\makeatletter
\@ifundefined{thm}{\newtheorem{thm}{Theorem}}{}
\@ifundefined{lem}{\newtheorem{lem}[thm]{Lemma}}{}
\makeatother
\newcommand{\bC}{\mathbb C}
\newcommand{\bZ}{\mathbb Z}
\newcommand{\cA}{\mathcal A}
\title[Relative monodromy of ramified sections on abelian schemes]{Relative monodromy of ramified sections on abelian schemes}
\author{Paolo Dolce, Francesco Tropeano}
\date{Source: arXiv:2407.19476v4 (CC BY 4.0)}
\begin{document}
\maketitle

\noindent\emph{Source and licence.} Relative monodromy of ramified sections on abelian schemes. Version arXiv:2407.19476v4 (CC BY 4.0), distributed under the Creative Commons Attribution 4.0 International licence, as recorded in the arXiv licence metadata for this exact version and shown on the abstract page https://arxiv.org/abs/2407.19476v4. Full rights notice: licenses/arxiv-2407.19476-CC-BY-4.0.txt.\par\smallskip

\noindent\textit{Editorial scope.} Counted unit: the Lem whose source label is \texttt{baseChange} in the article (source0.tex lines 795--797, source label \texttt{baseChange}), delivered together with the article's own complete original proof of it (source0.tex lines 799--830, one \texttt{proof} environment of 32 lines). No mathematics is written, repaired or re-derived: statement and proof are the article's own text line for line. Required context retained uncounted: mainThm1 (lines 219--221). Every definition, display and label of those lines is retained unchanged. Omitted from this package and not redistributed: 1--218, 222--794, 798--798, 831--1510. Nothing inside a retained excerpt is omitted or altered: every retained body is delivered exactly as the article's lines are written, so no omission receipt of any kind accompanies this package. Citations: the retained excerpts carry no citation command at all, so no bibliography entry of the article is transcribed and no reference is redistributed. Reference adaptation: none was needed and none was made; every reference inside the delivered unit resolves inside it, so no unresolved reference or citation is produced. Printed slips kept literal: no printed word, symbol, hypothesis or number of the counted statement or of its proof was altered. Layout and class: the article's own class is \texttt{article} with packages including appendix, amsmath, tocbibind, stmaryrd, verbatim, comment, mathrsfs, mathtools, amsthm, amssymb, mdframed, geometry, enumerate, bm; the wrapper is the lane's pinned \texttt{amsart} editorial wrapper at 10pt on A4, which restates the theorem-like environments the retained text needs and adds the article's own macro definitions that the retained text needs (\texttt{\textbackslash bC}, \texttt{\textbackslash bZ}, \texttt{\textbackslash cA}), with extra wrapper packages (bm, bbm, dsfont, xspace, cancel, mathtools, cleveref). The delivered numbering is the unit's own. Assets: no retained span contains a figure environment, a graphics-inclusion command, a PDF-page inclusion, a table or a TikZ picture; no image file is copied into this package, the package declares no asset dependency and no image file is redistributed with it. Licence: version arXiv:2407.19476v4 of this article is distributed under the Creative Commons Attribution 4.0 International licence, as recorded in the arXiv licence metadata for this exact version and shown on the abstract page https://arxiv.org/abs/2407.19476v4. A full rights notice is delivered at licenses/arxiv-2407.19476-CC-BY-4.0.txt. Characters: every retained excerpt is pure ASCII, so the delivered engine, XeLaTeX with its default Latin Modern text font, reports no missing character.
\begin{thm}\label{mainThm1}
    Let $\phi:\mathcal A \to S$ be an abelian scheme without fixed part such that the modular map $p:S \to T \subseteq \mathbb A_g$ is generically  finite. Let $U_p\subseteq S$ be the Zariski open subset on which $p$ is finite and flat. Then for any section $\sigma\colon S\to \mathcal A$  ramified over $U_p$, the relative monodromy group $M_\sigma^{\textrm{rel}}$ is non-trivial.
\end{thm}

\begin{lem}\label{baseChange}
	\Cref{mainThm1} is invariant by finite base change: in other words, if $\varphi: \widetilde{S} \rightarrow S$ is a finite morphism and $\widetilde{\cA} \to \widetilde{S}$ is the pullback of $\cA\rightarrow S$ via $\varphi$, then \Cref{mainThm1}  for $\widetilde{\cA} \rightarrow \widetilde{S}$ implies \Cref{mainThm1}  for $\cA \rightarrow S$.
\end{lem}

\begin{proof}
	Let us denote by $\omega_1, \ldots, \omega_{2g}$ a basis of periods of $\cA \rightarrow S$. We can define $\widetilde{\omega_1}, \ldots, \widetilde{\omega_{2g}}$, a basis of periods of $\widetilde{\cA} \rightarrow \widetilde{S}$, by the equations
	\begin{equation}\label{eqn3}
		\widetilde{\omega_1}=\omega_1 \circ \varphi,\; \ldots,\; \widetilde{\omega_{2g}}=\omega_{2g} \circ \varphi.
	\end{equation}
	Fix two base points $\widetilde{s} \in \widetilde{S}(\bC)$ and $s \in S(\bC)$ such that $\varphi(\widetilde{s})=s$. We have the associated monodromy representations:
	\[
    \rho:\pi_1(S(\bC),s) \rightarrow \textrm{Mon}(\cA), \qquad \widetilde{\rho}:\pi_1(\widetilde{S}(\bC),\widetilde{s}) \rightarrow \textrm{Mon}(\widetilde{A}).
    \]
	By \Cref{eqn3}, we have $\widetilde{\rho}(h)=\rho(\varphi_*(h))$ for each $h \in \pi_1(\widetilde{S}(\bC),\widetilde{s})$ where $\varphi_*:\pi_1(\widetilde{S}(\bC),\widetilde{s})\rightarrow \pi_1(S(\bC),s)$ denotes the induced homomorphism between fundamental groups. In particular, we obtain
	\begin{equation}\label{eqn3_bis}
		\varphi_*(\ker{\widetilde{\rho}})=\ker{\rho} \cap \varphi_*(\pi_1(\widetilde{S}(\bC),\widetilde{s})).
	\end{equation}
	Let $\sigma:S \rightarrow \cA$ be a non-torsion section of $\cA \rightarrow S$. Since $\widetilde{\cA}\rightarrow \widetilde{S}$ is obtained as pullback of $\cA \rightarrow S$, then the abelian varieties $\widetilde{\cA}_{\widetilde{s}}$ and $\cA_s$ are canonically identified; therefore the pullback $\varphi^*\sigma:\widetilde{S} \rightarrow \widetilde{\cA}$ is a non-torsion section of $\widetilde{\cA} \rightarrow \widetilde{S}$. We have the associated monodromy representations
	\[
    \theta_\sigma:\pi_1(S(\bC),s) \rightarrow \textrm{SL}_{2g+1}(\bZ), \qquad \theta_{\varphi^*\sigma}:\pi_1(\widetilde{S}(\bC),\widetilde{s}) \rightarrow \textrm{SL}_{2g+1}(\bZ).
    \]
		
	Using again the fact that the abelian varieties $\widetilde{\cA}_{\widetilde{s}}$ and $\cA_s$ are canonically identified, we obtain that a determination of the logarithm of $\varphi^*\sigma$ at $\widetilde{s}$ can be defined by the equation
	\begin{equation}\label{eqn4}
		\log_{\varphi^*\sigma}(\widetilde{s}):=\log_\sigma(s).
	\end{equation}
	Hence, we have $\theta_{\varphi^*\sigma}(h)=\theta_\sigma(\varphi_*(h))$ for each $h \in \pi_1(\widetilde{S}(\bC),\widetilde{s})$. Now, let us consider the relative monodromy group
	\[
	M_{\varphi^*\sigma}^{\textnormal{rel}}=\theta_{\varphi^*\sigma}(\ker{\widetilde{\rho}}).
    \]
    By \Cref{eqn3_bis} and \Cref{eqn4}, we obtain
	\[
	M_{\varphi^*\sigma}^{\textnormal{rel}} = \theta_\sigma(\varphi_*(\ker{\widetilde{\rho}})) = \theta_{\sigma}(\ker{\rho} \cap \varphi_*(\pi_1(\widetilde{S}(\bC),\widetilde{s}))).
    \]
    Since $\ker{\rho} \cap \varphi_*(\pi_1(\widetilde{S}(\bC),\widetilde{s})) \subset \ker{\rho}$, we get $M_{\varphi^*\sigma}^{\textnormal{rel}} \subseteq M_{\sigma}^{\textnormal{rel}}$. The conclusion follows.\\
\end{proof}

\end{document}
